\documentclass[../main.tex]{subfiles}
\begin{document}
\omk

\nk{Chỗ này\ldots đúng không nhỉ?}

Tôi đã đi loanh quanh khắp thành phố này gần ba tiếng đồng hồ, tính cả thời gian bắt xe buýt. Địa chỉ họ ghi loằng ngoằng quá! Mà\ldots đây có phải là nơi tôi sẽ làm sắp tới không nhỉ? 

Sau một ngày chuẩn bị, hôm nay tôi sẽ đến nơi này lần đầu tiên để nhận công việc. Tôi phải nắm rõ mọi thứ để không bị bất ngờ khi bước vào nơi làm việc của mình. Ngày mai sẽ là ngày đầu tiên tôi làm tại nơi này. Và tôi biết là sẽ đến nhận việc, nên tôi đã diện một cách lịch sự (thực ra là vì bố tôi bảo thế, nhưng ông ấy nói cũng đúng): một chiếc áo phông màu đỏ đi với một cái quần bò có cài thắt lưng và đi dép quai hậu. Tôi không thích đi giày, hôi chân lắm. Tóc tôi cũng được chải chuốt lại để ít nhất trông không quá lộn xộn.

Sau một hồi đi bộ, tôi chợt nhìn thấy một tòa nhà từ xa. Nó cao khoảng chín, mười tầng, giống như những tòa nhà khác, nhưng\ldots Chắc là đây rồi. Tòa nhà này có ghi rõ chữ \textit{hololive} to tướng, lại còn sơn màu trắng như trong bưu thiếp tôi nhận. Chắc chắn không thể nhầm lẫn được nữa.

\ct

Tôi mò mẫm bước vào trong, đến nơi mà các thành viên của hololive làm việc. Vậy nhưng khi vừa vào, tôi nhìn thấy ba cô gái đang ôm chầm lấy nhau, rõ ràng là rất run sợ. \nk{Cái gì vậy?} tôi thắc mắc. 

Vừa trông thấy tôi, cả ba người sợ hãi vội chạy tới. Họ là Minato Aqua, Natsuiro Matsuri và Roboco - ba cô gái mà tôi đã từng trao đổi qua mạng từ ngày hôm qua nhưng chưa có dịp gặp trực tiếp.

\begin{dialogue} 
    \speak{Aqua} \nk{Ơn giời anh đây rồi\ldots}
    \speak{Matsuri} \nk{Kuon-senpai, giúp bọn em với!} 
    \speak{Roboco} \nk{Bọn em vừa thấy một con gì đó bay vào phòng á!!} 
\end{dialogue}

Tôi nghiêng đầu ra, không hiểu họ đang nói gì. Tôi định bảo bọn họ giải thích lại mọi chuyện thì họ tiếp tục:

\begin{dialogue}
\speak{Roboco} Bọn em vừa thấy một con gì đó bay vào phòng á!!
\speak{Kuon} Tôi đớ người ra, chưa hiểu chuyện gì xảy ra. Tôi định hỏi họ thì Aqua tiếp lời:
\speak{Aqua} Nó to lắm á, trông như là con ong vậy!
\speak{Matsuri} Bọn em sợ nó đốt lắm!
\speak{Roboco} Anh làm gì với nó đi chứ!!
\end{dialogue}

Vừa dứt lời, thứ côn trùng gớm ghiếc xuất hiện ở phía sau tôi, cả ba người lập tức hét lên cảnh báo tôi: \nk{\textbf{Đằng sau anh kìa!}} Tôi quay lại nhìn và lập tức giơ cái dép lào phòng thủ từ trước (tại sao tôi lại mang đi nhỉ?) và\ldots vút! Chỉ trong giây lát, nó đã còng queo ở dưới sàn và chầu trời.

\begin{dialogue}
\speak{Kuon} Xời. Mấy con gián này ở nhà anh đầy. Dễ.
\speak{Aqua} N-Nó chết chưa vậy?
\end{dialogue}

Nhìn thấy thứ sinh vật kia không còn cựa quậy nữa, họ mới thở phào nhẹ nhõm.

\begin{dialogue}
    \speak{Matsuri} Được cứu rồi\ldots
    \speak{Roboco} Cảm ơn anh nhiều!
\end{dialogue}
    
Tôi chưng hửng nhìn Roboco và nói với giọng mỉa mai: \nk{Mà Roboco ạ, em là rô-bốt thì sợ quái gì bị đốt? Cái tia la-de trên tay để làm gì?}    

Ngay lập tức, cô nàng rô-bốt phản ứng lại: \nk{Anh muốn chúng em đền cả phòng à!?}

\nk{Ờ thì, cứ diệt đi là được. Đền bao nhiêu anh trả,} tôi đáp lại em ấy bằng một câu nói đùa. Tất nhiên, với độ PON của em ấy, Roboco không thể hiểu được: \nk{Thôi, em xin.} Cả Matsuri cũng gật đầu đồng tình với em ấy.
    
Nói vậy thôi, chứ tôi đào đâu ra tiền để đền chứ. Không phải tự dưng tôi đến đây xin việc làm đúng ngành mà lương cao đâu.

Sau khi chúng tôi giải quyết xong với con côn trùng đó, chúng tôi lên tới văn phòng và thấy Fubuki vẫn đang khám xét cái gì đó. 
\begin{dialogue}
\speak{Fubuki} A, mọi người đến rồi! Chào buổi sáng!
\end{dialogue}
Cả ba người họ đều đáp lại lời chào của nàng cáo trắng. 
Tôi chợt nhận ra điều gì đó, \nk{À, cái máy đó chắc là của A-san đấy. Chị ấy lắp ở đây để làm gì ta\ldots?} Tôi đi tới và cầm chiếc máy quay đó lên, nhưng nhìn nó đã tắt ngúm đi từ khi nào. 

\begin{dialogue}
\speak{Matsuri} Thật kỳ lạ, ai lại để máy quay ở đây chứ?
\speak{Aqua} Chắc là A-chan đặt nó ở đây để quay thử cảnh gì đó chăng? Nhưng sao nó lại tắt rồi?
\speak{Roboco} Em cũng không rõ nữa. A-chan là người hay làm mấy trò bất ngờ lắm mà, có khi là kế hoạch gì đó đấy.
\speak{Kuon} Hoặc là có thể hết pin. Anh đoán vậy. Nhưng vấn để là, {\color{blue}tại sao chứ?}
\speak{Fubuki} Đúng vậy á! Tại sao chị ấy lại để nó ở đây mà không nói với chúng ta gì hết nhỉ? Cứ như là\ldots muốn chúng ta tự tìm hiểu vậy.
\speak{Matsuri} Cái này nghe có vẻ giống một trò chơi ghê á.
\speak{Aqua} Anh nghĩ sao về cái máy này, Kuon-senpai?
\speak{Kuon} Anh nghĩ\ldots chúng ta không nên động vào nó. Cứ để A-san giải thích sau.
\end{dialogue}

Nói rồi, tôi cất chiếc máy quay đi và hỏi bốn người về dự án tiếp theo của họ.
\begin{dialogue}
\speak{Kuon} Vậy dự án\ldots cái gì, "Holo no Graffiti" triển khai đến đâu rồi? Anh có nghe Tanigo-san nói thế sáng nay\ldots
\end{dialogue}

Cả bốn người đồng loạt nhìn nhau, rồi quay lại nhìn tôi với vẻ mặt ngơ ngác.
\begin{dialogue}
\speak{Matsuri} holo no Graffiti? Chưa nghe tên này bao giờ\ldots
\speak{Aqua} Là cái gì vậy Kuon-senpai? Dự án mới à?
\speak{Roboco} Chắc chắn không phải cái gì mà chúng ta đã làm rồi đúng không?
\speak{Fubuki} Chúng ta có biết cái tên này không nhỉ? Mình cũng không nhớ đã nghe qua.
\end{dialogue}

Tôi hơi ngạc nhiên vì các cô gái không biết gì về dự án này. Hóa ra cái tên \nk{holo no Graffiti} là một cái gì đó hoàn toàn mới đối với họ. Mà thật ra, tôi cũng chẳng hiểu rõ lắm về nó, chỉ nghe qua từ cuộc điện thoại sáng nay với A-san và Tanigo-san. Hoặc là họ giải thích khá qua loa, hoặc là do tôi đang vội.

\begin{dialogue}
\speak{Kuon} Thật sự anh cũng chẳng rõ lắm về dự án này\ldots A-chan và Tanigo chỉ nói sơ qua trong cuộc gọi sáng nay. Hình như là một cái gì đó liên quan đến nghệ thuật hay video gì đó, nhưng\ldots anh cũng không rõ lắm.
\speak{Matsuri} Ối chà\~{}, thế thì anh phải giải thích cho chúng em nghe đi chứ!
\speak{Aqua} Đúng đó! Anh biết gì về nó thì chia sẻ đi!
\speak{Roboco} Em muốn biết thêm nữa, nghe cứ như là một dự án lớn ấy.
\end{dialogue}

Tôi thấy hơi bối rối, nhưng vẫn cố gắng trả lời.

\begin{dialogue}
\speak{Kuon} Để tôi thử xem sao\ldots Nhưng anh chỉ biết tên dự án và rằng nó sắp được triển khai thôi. Hình như là bắt đầu từ tuần sau. Hôm nay mới chỉ là demo thôi.
\end{dialogue}

\nk{Vậy à\ldots} bọn họ nhìn nhau với những con mắt tò mò. \nk{Vậy chúng ta triển thôi chứ?}

Cả bốn người đều trả lời đồng thanh:
\begin{dialogue}
\speak{4 người} Hai\~{}! (Vâng\~{}!)
\speak{Kuon} Tốt. Vậy thì\ldots ta khởi hành nào. (Let's roll.)
\end{dialogue}

Và rồi, cuộc phiêu lưu của tôi - Hikawa Kuon và các thành viên hololive chính thức bắt đầu. 

\ldots Ít nhất là để tôi nhận chính thức công việc từ CEO Tanigo-san trước đã.

\ct

Tôi tìm đến văn phòng CEO trong tòa nhà, được chỉ dẫn bởi A-chan đang làm ở đây. Tại đó, tôi và Tanigo cùng chị ấy bàn nhau về lịch thời gian biểu của tôi làm việc ở công ty. 

Tôi làm bán thời gian từ thứ hai đến thứ sáu, buổi sáng từ tám giờ đến mười một giờ hai mươi lăm. Buổi chiều từ một giờ đến sáu giờ. Nhà tôi cách công ty bằng một chuyến xe buýt dài bốn mươi phút, cộng thêm mười phút đi bộ, nên tôi thường sẽ có mặt tại đó trước tận\ldots một tiếng, đấy là nếu tôi không có tiết học.

Với hai ngày cuối tuần, tôi làm toàn thời gian, từ tám giờ ba mươi sáng đến năm giờ chiều. Tất nhiên, tôi sẽ cố gắng để kéo dài thời gian đến sáu giờ chiều mới được đi về, các bạn có thể gọi tôi là kẻ cuồng công việc cũng không sai.

YAGOO và A-san cũng hỏi tôi một chút về bản thân. Tôi có nói rằng điểm mạnh của tôi là trung thực, thẳng thắn (dù hơi quá), và cần cù làm việc. Tôi cũng tự nhận có trí thông minh khá tốt, và cũng đã cho hai người xem bằng cách lập trình. Tôi cũng có thể nói được tiếng Anh (thực tế tôi có thể nói được nhiều thứ tiếng hơn, nhưng tôi giấu nhẹm đi vì có thể tôi sẽ không dùng đến\ldots đúng không nhỉ?)

Như vậy, giấy tờ cũng đã ký xong, công việc cũng đã được điều sẵn, tôi chính thức trở thành tổng quản lý - một công việc mà tôi không hề biết gì, nhưng lại là công việc đáng mơ ước của đám con trai, điều mà tôi cũng không biết - đồng hành với Yuujin A-san.

Còn về vụ \nk{holo no Graffiti} ấy thì sao? Tôi đã hỏi cả ông CEO lẫn đồng nghiệp mới của tôi, nhưng họ chỉ nói với tôi rằng: \nk{Cậu hãy tự mình trải nghiệm đi rồi sẽ biết.}

Giờ tôi chỉ có thể tự chúc phúc bản thân để có thể\ldots ờm, xứng đáng với chức vụ này và\ldots mang tiền về cho bố mẹ\ldots? Chắc thế.

Mà thật lòng thì, tôi đã tự đưa mình vào tình huống quái thai nào thế này?

\end{document}